\chapter{Research Hypotheses, Questions \& Pre-Phase A
Closure}\label{chapter-5-research-hypotheses-questions--pre-phase-a-closure}

\begin{quote}
\textbf{Writing Guide \& Methodological Focus:}

\begin{itemize}
\item
  \textbf{Objective:} To formally state the central hypothesis and the
  specific, testable questions that the research will answer. State the
  science that will be tested, then close \textbf{Pre-Phase A} with a
  light, systems-engineering frame that shows the methods that will be
  used.
\item
  \textbf{Method:} Short synthesis → one-line hypothesis → concise
  research questions → hint-level Pre-Phase-A closure
  (needs/goals/objectives → constraints → success criteria) → note the
  concept options to be tested.
\item
  \textbf{Caution:} This chapter should be short and direct. It is the
  pivot point between the background and this thesis's original work.
  Keep it brief and directional; traceable requirements IDs, success
  criteria and tables here.
\item
  \textbf{Checkpoint:} Is the methodology described in Part II directly
  answerable to the research questions? After this chapter, the reader
  knows what will be tested* and \emph{the frame it must satisfy}.
\end{itemize}
\end{quote}

\section{Synthesis of the
Problem}\label{51-synthesis-of-the-problem}

Update the Chapter 1 Problem Statement with condensed information from
Chapters 2--4 into a single Problem Statement paragraph

\subsection{Primary Research
Hypothesis}\label{52-primary-research-hypothesis}

Assert that a \textbf{Portugal-tuned, tip-and-cue LEO all-weather
wildfire constellation} can meet the need more effectively than
conventional single-layer designs.

\section{Key Research
Questions}\label{53-key-research-questions}

List of questions; cadence vs orbit geometry; constellation sizing/plane
spacing; single vs two-layer (tipper+cuer); angle/quality policy impact;
robustness under typical cloud/smoke; latency pathway; transferability
to secondary ROIs.

\section{\texorpdfstring{Pre-Phase-A Closure
}{Pre-Phase-A Closure }}\label{54-pre-phase-a-closure}

Clarify that we are using the Nasa System Engineering approach to
develop the concept of this mission. Explain the evolution of this
structure. Add the system engineering thesis structure.

\subsection{Needs, goals,
objectives}\label{541-needs-goals-objectives}

Freeze a brief structure condensing all mission requirements build into
here

\subsection{Constraints}\label{542-constraints}

Note the boundary conditions you will work

\subsection{Success criteria}\label{543-success-criteria}

Table with success for each requirement

\subsection{Mission Concept
Options}\label{544-mission-concept-options}

From the background development, suggest concepts to be verified in the
development of Part II: Methodology and Mission Design, such as
constellation type (single vs. two-layer), orbit class, plane/phasing
patterns, and high-level sensor families/limits.
